% ═══════════════════════════════════════════════════════════════
%  ЦВЕТОВАЯ СХЕМА
% ═══════════════════════════════════════════════════════════════


% ─── Семантические цвета ──────────────────────────────────
\definecolor{clrText}{HTML}{1A1A1A}        % основной текст
\definecolor{clrCodeBg}{HTML}{F5F5F5}      % фон кода
\definecolor{clrCodeFrame}{HTML}{CCCCCC}   % рамка кода
\definecolor{clrKeyword}{HTML}{0055AA}     % ключевые слова
\definecolor{clrComment}{HTML}{228855}     % комментарии
\definecolor{clrString}{HTML}{AA2222}      % строки

\definecolor{clrIdea}{HTML}{2E75B6}        % идея (синий)
\definecolor{clrWarn}{HTML}{C0392B}        % внимание (красный)
\definecolor{clrKey}{HTML}{27AE60}         % ключевой момент (зелёный)
\definecolor{clrNote}{HTML}{8E44AD}        % заметка (фиолетовый)

\definecolor{clrUIData}{HTML}{2980B9}      % UI (синий)
\definecolor{clrLogicData}{HTML}{27AE60}   % логика (зелёный)
\definecolor{clrHWData}{HTML}{E67E22}      % железо (оранжевый)
\definecolor{clrStateData}{HTML}{8E44AD}   % состояние (фиолетовый)

\definecolor{clrFlowData}{HTML}{C0392B}    % поток данных (красный)
\definecolor{clrFlowCmd}{HTML}{2980B9}     % поток команд (синий)
\definecolor{clrDep}{HTML}{7F8C8D}         % зависимости (серый)

% Применяем основной цвет текста
\color{clrText}

% ═══════════════════════════════════════════════════════════════
%  ЗАГОЛОВКИ
% ═══════════════════════════════════════════════════════════════
\usepackage{titlesec}

% Глава
\titleformat{\chapter}[display]
    {\normalfont\huge\bfseries\color{clrIdea}}
    {\filleft\Large\MakeUppercase{\chaptertitlename}~\thechapter}
    {1ex}
    {\titlerule\vspace{1ex}\filright}
    [\vspace{1ex}\titlerule]

\titlespacing*{\chapter}{0pt}{-20pt}{30pt}

% Раздел
\titleformat{\section}
    {\normalfont\Large\bfseries\color{clrIdea}}
    {\thesection}{1em}{}

\titleformat{\subsection}
    {\normalfont\large\bfseries\color{clrText}}
    {\thesubsection}{1em}{}

\titleformat{\subsubsection}
    {\normalfont\normalsize\bfseries\color{clrText}}
    {\thesubsubsection}{1em}{}
    

% ─── Общие стили узлов (для переиспользования) ────────────
\tikzset{
    % Базовый модуль
    module/.style={
        draw, rounded corners=1.5mm, thick,
        minimum width=2.5cm, minimum height=0.9cm,
        align=center, font=\footnotesize\ttfamily,
        fill=white,
        drop shadow={opacity=0.15, shadow xshift=0.5mm, shadow yshift=-0.5mm},
    },
    % Модули по слоям
    modUI/.style={module, fill=clrUIData!20},
    modLogic/.style={module, fill=clrLogicData!20},
    modHW/.style={module, fill=clrHWData!25},
    modState/.style={module, fill=clrStateData!20},
    modUtil/.style={module, fill=gray!15},
    modEntry/.style={module, fill=clrWarn!20},
    % Внешние зависимости
    ext/.style={
        draw, dashed, rounded corners=1.5mm,
        minimum width=2cm, minimum height=0.8cm,
        align=center, font=\footnotesize\ttfamily,
        fill=clrUIData!5,
    },
    % Стрелки
    dep/.style={-{Latex[length=2mm]}, thick, color=clrDep},
    extdep/.style={dep, dashed, color=clrUIData!60},
    flowData/.style={-{Latex[length=3mm]}, very thick, color=clrFlowData},
    flowCmd/.style={-{Latex[length=3mm]}, very thick, color=clrFlowCmd, dashed},
    % Метки
    label/.style={font=\footnotesize, midway, fill=white, inner sep=1pt},
}

% ═══════════════════════════════════════════════════════════════
%  СТИЛИ ДЛЯ ДИАГРАММ ПОСЛЕДОВАТЕЛЬНОСТИ (UML sequence)
% ═══════════════════════════════════════════════════════════════
\tikzset{
    % ─── Участник (актор) ────────────────────────────────
    seqactor/.style={
        draw, rounded corners=1.5mm, thick,
        minimum width=2.4cm, minimum height=0.7cm,
        align=center, font=\footnotesize\ttfamily,
        fill=white,
    },
    seqactoruser/.style={seqactor, fill=clrWarn!20},
    seqactorkernel/.style={seqactor, fill=clrUIData!20},
    seqactorgio/.style={seqactor, fill=clrStateData!20},
    seqactorsp/.style={seqactor, fill=clrHWData!25},
    seqactorcb/.style={seqactor, fill=clrLogicData!25},
    % ─── Линия жизни ────────────────────────────────────
    lifeline/.style={
        thick, dashed, color=gray!60,
    },
    % ─── Сообщение (стрелка) ─────────────────────────────
    msg/.style={
        -{Latex[length=2.5mm]}, thick, color=black!80,
    },
    msguser/.style={msg, color=clrWarn!80!black},
    msgkernel/.style={msg, color=clrUIData!80!black},
    msggio/.style={msg, color=clrStateData!80!black},
    msgsp/.style={msg, color=clrHWData!80!black},
    msgcb/.style={msg, color=clrLogicData!80!black},
    % ─── Возврат (пунктир) ───────────────────────────────
    ret/.style={
        -{Latex[length=2mm]}, thick, dashed, color=gray!70,
    },
    % ─── Метка сообщения ────────────────────────────────
    msglabel/.style={
        font=\scriptsize\ttfamily,
        fill=white, inner sep=1pt,
        above,
    },
    msglabelbelow/.style={
        font=\scriptsize\ttfamily,
        fill=white, inner sep=1pt,
        below,
    },
    % ─── Блок активации (прямоугольник на линии) ────────
    activation/.style={
        fill=white, draw, thick, minimum width=0.25cm,
    },
}

% ═══════════════════════════════════════════════════════════════
%  ВРЕЗКИ
% ═══════════════════════════════════════════════════════════════
\usepackage[most]{tcolorbox}
\tcbuselibrary{skins, breakable}

% ─── Идея ─────────────────────────────────────────────────
\newtcolorbox{ideabox}[1][]{
    enhanced,
    colback=clrIdea!5,
    colframe=clrIdea,
    colbacktitle=clrIdea,
    coltitle=white,
    fonttitle=\bfseries,
    title=\faLightbulb~Идея,
    boxrule=0.5pt,
    arc=2mm,
    left=2mm, right=2mm, top=2mm, bottom=2mm,
    breakable,
    #1
}

% ─── Внимание ─────────────────────────────────────────────
\newtcolorbox{warnbox}[1][]{
    enhanced,
    colback=clrWarn!5,
    colframe=clrWarn,
    colbacktitle=clrWarn,
    coltitle=white,
    fonttitle=\bfseries,
    title=\faExclamationTriangle~Внимание,
    boxrule=0.5pt,
    arc=2mm,
    left=2mm, right=2mm, top=2mm, bottom=2mm,
    breakable,
    #1
}

% ─── Ключевой момент ──────────────────────────────────────
\newtcolorbox{keybox}[1][]{
    enhanced,
    colback=clrKey!5,
    colframe=clrKey,
    colbacktitle=clrKey,
    coltitle=white,
    fonttitle=\bfseries,
    title=\faKey~Ключевой момент,
    boxrule=0.5pt,
    arc=2mm,
    left=2mm, right=2mm, top=2mm, bottom=2mm,
    breakable,
    #1
}

% ─── Заметка ──────────────────────────────────────────────
\newtcolorbox{notebox}[1][]{
    enhanced,
    colback=clrNote!5,
    colframe=clrNote,
    colbacktitle=clrNote,
    coltitle=white,
    fonttitle=\bfseries,
    title=\faInfoCircle~Заметка,
    boxrule=0.5pt,
    arc=2mm,
    left=2mm, right=2mm, top=2mm, bottom=2mm,
    breakable,
    #1
}

% ─── Упражнение ───────────────────────────────────────────
\newtcolorbox{exercisebox}[1][]{
    enhanced,
    colback=clrKey!3,
    colframe=clrKey!70!black,
    colbacktitle=clrKey!70!black,
    coltitle=white,
    fonttitle=\bfseries,
    title=\faPen~Упражнение,
    boxrule=0.5pt,
    arc=2mm,
    left=2mm, right=2mm, top=2mm, bottom=2mm,
    breakable,
    #1
}

% ═══════════════════════════════════════════════════════════════
%  ТАБЛИЦЫ
% ═══════════════════════════════════════════════════════════════
\usepackage{booktabs}
\usepackage{tabularx}
\usepackage{longtable}
\usepackage{array}
\usepackage{multirow}
\usepackage{colortbl}

% ─── Цветной заголовок таблицы ────────────────────────────
\newcommand{\tablehead}[1]{\textcolor{white}{\textbf{#1}}}

% ─── Стиль «грабли» ───────────────────────────────────────
\newcolumntype{L}{>{\raggedright\arraybackslash}X}
\newcolumntype{C}{>{\centering\arraybackslash}X}
\newcolumntype{R}{>{\raggedleft\arraybackslash}X}

% ─── Шапка таблицы с фоном ────────────────────────────────
\newcommand{\headrow}[1]{%
    \rowcolor{clrIdea!80}%
    \tablehead{#1}%
}

%  Эти стили дают нам:
%- 5 цветных акторов — каждый со своим цветом.
%- Линии жизни — серые пунктирные.
%- Сообщения — сплошные стрелки, цветные по отправителю.
%- Возвраты — серые пунктирные стрелки.
%- Метки — моноширинные, с белым фоном.
